\documentclass[11pt]{article}
\usepackage[margin=1in]{geometry}
\usepackage{amsmath,amssymb}
\usepackage{booktabs,array}
\usepackage{xurl}
\usepackage[hidelinks]{hyperref}
\setlength{\emergencystretch}{2em}
\title{A low-speed regime of the vehicle model: the vehicle can stop}
\author{}
\date{October 11, 2026}
\begin{document}
\maketitle
\paragraph{Status and scope.}
This is the second step approved by the user on October 11, 2026, after the
mild terminal set of commit \texttt{8d43230}. The terminal set treats a mild
state as kinematic, and it contains stops: the ego brakes to rest behind a
vehicle that stays in the range. The declared model, and with it the
simulated plant, ended at 1~m/s and could not reach those states. This step
gives the declared model a kinematic regime at low speed. It reports the
checks, a probe in which stopping is the only safe outcome, and a paired
campaign against \texttt{8d43230}.

\paragraph{Result.}
On a road too narrow to pass, behind a stationary vehicle 30~m ahead at
8~m/s, the controller now stops 0.23~m behind it and stays there for the
30~s of the probe. Before this step it failed at 1.07~m/s after 3.8~s. The
campaign is nearly unchanged: 13 of 14 exact encounters (the same ones) and
72 of 84 noisy (73 before), with no collision. The campaign seldom slows
the vehicle below 4~m/s: 11 encounters dip to 3.6--3.9~m/s, and none goes
lower. The differences come from numerical failures of the conic programs and a
different plan in one exact crossing, whose least replay clearance is
0.001~m (0.140~m before) at speeds above 6~m/s.

\section{The model}
The declared model (\texttt{nonlinearBicycleModel.derivative}) is the Fiala
bicycle above \texttt{model.kinematicSpeedRange}~$=[2.5,4]$~m/s and a
kinematic bicycle below it. Between the two it blends with the cubic
smoothstep weight $\lambda(v_x)$. Position and heading follow the body
velocities in all three. Above 4~m/s it is the Fiala bicycle exactly, so the
CLF certificates at 8 and 15~m/s are unaffected. The Fiala tire's own floor
(\texttt{scheduleSpeedFloor}, 1~m/s) lies inside the kinematic regime, where
the Fiala forces are not evaluated.

The kinematic bicycle has three parts.
\begin{itemize}
\item \emph{Lateral states.} The yaw rate relaxes with
  \texttt{kinematicLagSeconds}~$=0.1$~s to $v_x\tan\delta/L$. The lateral
  velocity relaxes to $l_r$ times that value, which means no slip at the
  rear axle.
\item \emph{Longitudinal force.} It is the braking ratio times the axles'
  force scales (the front one through $\cos\delta$), with the road load and
  $v_yr$.
\item \emph{Braking.} It fades as $\tanh(v_x/0.25)$
  (\texttt{brakeHoldSpeed}), so it ends at rest and never reverses the
  vehicle; driving does not fade, so the vehicle can move off from rest.
\end{itemize}

\paragraph{Exact rest.}
A vehicle that is braked and slower than 5~cm/s at the end of a hold is at
rest exactly: its body velocities and their sensitivities are set to zero.
Without this step, the fading brake leaves a tail of velocities and
coefficients of 1e-4 to 1e-7. In the probe below that tail left the CLF
stage's cone program without a numerical result, about 1~s after the
vehicle had stopped, and again while it crept to rest under light braking.
A diagnostic run that snapped the plant to rest kept the solves at
residuals of $5\cdot10^{-13}$, and with exact rest in the model the probe
holds the stop to its end. The
simulated plant (\texttt{ode45} on the same derivative) has no such step and
creeps the tail, about 1.5~cm under light braking.

\paragraph{Rows.}
The state's speed bound is now \texttt{speedMinimum} (0) instead of the
Fiala floor. In the kinematic regime the rear-adhesion cone reduces to
$|b|\le1$, because the rear tire has no lateral force there to saturate. The
sideslip cone applies above 0.25~m/s, because at rest the sideslip is
undefined, and the terminal set's mildness check skips it there as well.

\paragraph{Compiled kernel.}
The kernel's parameter vector gains the four constants (23 entries). The
interpreted model checks a compiled kernel against itself at a probe in each
regime, and a kernel that fails the call (another parameter layout) is not
used. An older checkout, whose check has no such guard, needs its kernels
rebuilt (\texttt{scripts/buildControllerKernels.m}).

\section{Checks}
\begin{itemize}
\item Full suite: 612 of 612 (\path{tests.txt}; 608 before). The new tests
  are listed below.
  \begin{itemize}
  \item The analytic tangent against finite differences in the blended
    (3.2~m/s) and kinematic (1.5 and 0.1~m/s) regimes, braking and
    driving, on straight and curved charts.
  \item Braking at the terminal law's ratio from 6~m/s to rest: never in
    reverse, the handling rows met, then exactly at rest with zero
    sensitivities. The distance is 4.4~m against 4.5~m for constant
    braking at $a_b$, because road load adds to the braking and the brake
    fades.
  \item At rest the vehicle stays still under braking, and moves off under
    driving.
  \item Above the kinematic range the derivative does not depend on the
    range.
  \end{itemize}
  The kernel test now draws speeds from 0 to 14~m/s.
\item Offline: the worst relative tangent error is $4.6\cdot10^{-10}$. From
  6~m/s the vehicle is below 1~mm/s after 1.8~s (1.5~s at a constant
  4~m/s$^2$).
\end{itemize}

\section{Stopping behind a stationary vehicle}
The probe uses a straight road with 1.2~m of lateral clearance on each side
of the lane centre, too narrow to pass. A stationary vehicle stands in the
lane 30~m ahead, and the ego starts at 8~m/s. The plant is the declared model
with RK4 holds, and the run lasts 600 frames.
\begin{itemize}
\item \texttt{8d43230}: the plan ran out after 77 frames (3.8~s) at 1.07~m/s
  and 0.38~m from the vehicle (\texttt{anchorUnavailable}).
\item This step, before exact rest: the ego stopped 0.23~m behind the
  vehicle after about 5~s, and the CLF stage failed numerically about 1~s
  later.
\item This step: the ego stops 0.23~m behind the vehicle and stays at rest
  to frame 600. The least clearance is 0.20~m.
\end{itemize}
The same probe at 15~m/s with the vehicle 45~m ahead fails at the first
frame in both versions (\texttt{pcbfNoNumericalResult}). Its startup
rollout cannot pass on this road, and this step does not address that.

\section{Campaign}
The 98 encounters were run with \texttt{8d43230} (``mild'') and with this
step (``low speed''). The base arm used kernels built from its own source.
The new arm was run twice: before and after exact rest. Both runs gave the
same outcome in every encounter; the second is reported.

\begin{table}[h]
\centering\small
\caption{Paired campaign. No collision in either arm.}
\begin{tabular}{lrrrr}
\toprule
 & \multicolumn{2}{c}{exact (14)} & \multicolumn{2}{c}{noisy (84)}\\
 & mild & low speed & mild & low speed\\
\midrule
recovered & 13 & 13 & 73 & 72\\
no solution: \texttt{nonlinearAcceptanceFailed} & 0 & 0 & 2 & 1\\
no solution: \texttt{anchorUnavailable} & 1 & 1 & 0 & 0\\
no solution: \texttt{pcbfNoNumericalResult} & 0 & 0 & 8 & 9\\
no solution: \texttt{clfNoNumericalResult} & 0 & 0 & 1 & 2\\
least replay clearance (m) & 0.140 & 0.001 & 0.391 & 0.391\\
median per-run least clearance (m) & 0.455 & 0.457 & 1.849 & 1.843\\
horizon, 95th percentile holds & 45 & 45 & 79 & 89\\
controller time per frame, 95th percentile (s) & 0.046 & 0.049 & 0.122 & 0.115\\
\bottomrule
\end{tabular}
\end{table}

Per geometry, the exact results are identical. The noisy results differ in
two geometries. Accelerating head-ons go from 11 to 12 of 12. Braking leads
go from 11 to 9 of 12: both 8 and 15~m/s with seed 3 now fail on
\texttt{clfNoNumericalResult}, the 15-m/s one at the first frame, and the
8-m/s one while following at 3.9~m/s in the blended regime.
\path{comparison.txt}, \path{campaign-runs.csv} and \path{failures.txt} list
every encounter.

\paragraph{The exact crossing at 0.001 m.}
In the exact 8-m/s crossing the vehicle never drops below 5.96~m/s, so it
never enters the low-speed regime. The plans differ nonetheless, because a
plan's braking tail now runs into the kinematic regime. From 1.2~s on, the
new arm chose a sharper swerve (steering to $-0.73$~rad), and the replay
clearance dropped to 0.9~mm, against 0.14~m before. Both are below the
0.2-m margin, which binds only at the plan's sampled nodes, and in the prefix
only as a minimized slack. No collision occurred.

\section{Decisions for the user}
\begin{itemize}
\item Keep this model as the declared model and the plant (committed). It is
  what lets a stop behind a vehicle be reached at all.
\item Below 4~m/s the estimator still assumes the Fiala force balance, and
  its domain ends at 1~m/s. No noisy encounter slowed below 3.5~m/s, so
  estimator-in-the-loop stops are untested.
\item Numerical failures of the two conic stages remain the main loss (12
  of 84 noisy encounters). Curved head-ons alone account for 9, and they
  predate this step.
\end{itemize}

\section{Reproduction}
\begin{itemize}
\item Kernels: \texttt{matlab -batch "addpath('scripts'); buildControllerKernels"}.
\item Unit tests: \texttt{matlab -batch "results = runtests('tests'); assertSuccess(results)"}.
\item Campaign: \path{scripts/run-workers.sh 2000 3} of this folder (base arm
  a \texttt{git archive} of \texttt{8d43230} with its own kernels, new arm a
  copy of the working tree), then
  \texttt{python3 -I scripts/compare.py <out-base> <out-new> <folder> mild lowspeed}.
  Raw outputs under \path{~/.cache/collisionAvoidance/low_speed_20261011/}
  are disposable and were deleted after the summaries were copied here.
\item The stopping probe was a session script and is not kept; its setup is
  described above.
\end{itemize}
\end{document}
